\documentclass[reqno, 12pt]{amsart}
\usepackage[brazil]{babel}
\usepackage[utf8]{inputenc}
\usepackage[T1]{fontenc}
\usepackage{lmodern}
\usepackage{amsfonts,amssymb,amsmath,textcomp,amsthm}
\usepackage{tikz}
\usetikzlibrary{calc}
\usepackage{geometry}
\geometry{left=25mm, right=25mm, top=22mm, bottom=28mm}
\usepackage{setspace}
\onehalfspacing
\usepackage{csquotes}
\usepackage[hidelinks]{hyperref}
\usepackage[backend=biber, style=numeric, maxbibnames=99]{biblatex}
\addbibresource{refs_antiramsey.bib}
\setlength{\parindent}{1.2cm}
\numberwithin{equation}{section}

% -------------------------------------------------------
% Ambientes
% -------------------------------------------------------
\theoremstyle{plain}
\newtheorem{theorem}{Teorema}[section]
\newtheorem{proposition}[theorem]{Proposição}
\newtheorem{lemma}[theorem]{Lema}
\newtheorem{corollary}[theorem]{Corolário}

\theoremstyle{definition}
\newtheorem{definition}[theorem]{Definição}
\newtheorem{question}[theorem]{Questão}
\newtheorem{problem}[theorem]{Problema}
\newtheorem{remark}[theorem]{Observação}
\newtheorem{example}[theorem]{Exemplo}

% -------------------------------------------------------
% Macros
% -------------------------------------------------------
\newcommand{\Gnp}{\mathcal{G}(n,p)}
\newcommand{\rb}{\xrightarrow{\,\mathrm{rb}\,}}
\newcommand{\ramsey}[1]{\xrightarrow{\,#1\,}}
\newcommand{\cram}{\xrightarrow{\,\mathrm{c\text{-}ram}\,}}
\newcommand{\mdens}{m_2}
\newcommand{\maxdens}{m}
\newcommand{\whp}{a.a.s.}

\date{\today}

% ============================================================
\begin{document}
% ============================================================

\begin{center}
  {\large\textbf{Ramsey Assimétrico e Anti-Ramsey em Grafos Aleatórios}}

  \vspace{0.4cm}

  Afonso Sant'Anna

  \vspace{0.2cm}

  {\small Orientador: Yoshiharu Kohayakawa \quad
   Coorientador: Guilherme O.\ Mota}

  \vspace{0.2cm}

  {\small Instituto de Matemática e Estatística -- USP}
\end{center}

\vspace{0.5cm}

\noindent\textbf{Nota.}
Este documento é um rascunho de trabalho para o doutorado de Afonso Sant'Anna.
Ele reúne definições, teoremas e resultados dos artigos da literatura relevante,
organizados para facilitar a escrita dos próprios resultados do aluno.
Seções marcadas com \emph{(RA)} referem-se ao tema de Ramsey Assimétrico/Canônico;
seções marcadas com \emph{(AR)} referem-se ao tema Anti-Ramsey.

\tableofcontents

\bigskip

% ============================================================
\section{Preliminares e Notação}
% ============================================================

Denotamos por $\Gnp$ o grafo aleatório binomial com $n$ vértices em que cada aresta
está presente independentemente com probabilidade $p = p(n)$.
Dizemos que $\Gnp$ satisfaz uma propriedade \emph{assintoticamente quase certamente}
(\whp) se a probabilidade de satisfazê-la tende a $1$ quando $n \to \infty$.

\begin{definition}[Limiar e tipos de limiar]
Seja $\mathcal{P}$ uma propriedade monótona crescente de grafos.
Uma função $\hat{p} = \hat{p}(n)$ é um \textbf{limiar grosseiro} (\emph{coarse threshold})
para $\mathcal{P}$ em $\Gnp$ se
\[
  \lim_{n\to\infty} \Pr[\Gnp \in \mathcal{P}] =
  \begin{cases}
    0 & \text{se } p = o(\hat{p}),\\
    1 & \text{se } p = \omega(\hat{p}).
  \end{cases}
\]
Dizemos que $\hat{p}$ é um \textbf{limiar semi-agudo} (\emph{semi-sharp threshold})
se existem constantes $c, C > 0$ tais que
\[
  \lim_{n\to\infty} \Pr[\Gnp \in \mathcal{P}] =
  \begin{cases}
    0 & \text{se } p \le c\,\hat{p},\\
    1 & \text{se } p \ge C\,\hat{p}.
  \end{cases}
\]
O \textbf{$0$-enunciado} refere-se à afirmação para $p$ abaixo do limiar,
e o \textbf{$1$-enunciado} à afirmação para $p$ acima do limiar.
\end{definition}

\begin{definition}[Densidades de um grafo]
\label{def:densidades}
Dado um grafo $H$, definem-se:
\begin{itemize}
  \item A \textbf{densidade} de $H$:
    $d(H) := \dfrac{e(H)}{v(H)}$, onde $e(H) = |E(H)|$ e $v(H) = |V(H)|$.

  \item A \textbf{densidade máxima} de $H$:
    \[
      \maxdens(H) := \max\!\left\{\frac{e(J)}{v(J)} : J \subseteq H,\; v(J) \ge 1\right\}.
    \]
    Dizemos que $H$ é \textbf{balanceado} se $\maxdens(H) = d(H)$, e
    \textbf{estritamente balanceado} se $\maxdens(H) > d(J)$ para todo
    $J \subsetneq H$ com $v(J) \ge 1$.

  \item A \textbf{2-densidade} de $H$ (para $v(H) \ge 3$):
    \[
      d_2(H) := \frac{e(H)-1}{v(H)-2}.
    \]

  \item A \textbf{2-densidade máxima} de $H$:
    \[
      \mdens(H) := \max\!\left\{d_2(J) : J \subseteq H,\; v(J) \ge 3\right\}
        \cup \left\{\tfrac{1}{2}\right\}.
    \]
    (A fração $\frac{1}{2}$ corresponde ao caso $H \cong K_2$.)
    Dizemos que $H$ é \textbf{$2$-balanceado} se $\mdens(H) = d_2(H)$, e
    \textbf{estritamente $2$-balanceado} se $\mdens(H) > d_2(J)$ para todo
    $J \subsetneq H$.
\end{itemize}
\end{definition}

\begin{example}
  Para o clique $K_k$ com $k \ge 3$:
  $\mdens(K_k) = \dfrac{\binom{k}{2}-1}{k-2} = \dfrac{k+1}{2}$.
  Para o ciclo $C_\ell$ com $\ell \ge 3$:
  $\mdens(C_\ell) = \dfrac{\ell-1}{\ell-2}$.
  Em particular,
  $\mdens(K_5) = 3$, $\mdens(C_5) = \frac{4}{3}$, $\mdens(C_4) = \frac{3}{2}$.
\end{example}

% ============================================================
\section{Propriedade de Ramsey em Grafos Aleatórios}
% ============================================================

\begin{definition}[Propriedade $r$-Ramsey]
Para grafos $G$ e $H$ e um inteiro $r \ge 2$, escrevemos
$G \ramsey{r} H$ para denotar a propriedade de que,
em toda coloração das arestas de $G$ com $r$ cores,
existe pelo menos uma cópia monocromática de $H$ em $G$.
\end{definition}

O seguinte resultado é o marco fundamental da teoria:

\begin{theorem}[Rödl--Ruciński \cite{rodl1993lower, rodl1994random, rodl1995threshold}]
\label{thm:rodl-rucinski}
Seja $r \ge 2$ um inteiro e seja $H$ um grafo que não é uma floresta de estrelas
(e, no caso $r = 2$, também não é um caminho de comprimento~$3$).
Existem constantes positivas $c = c(H,r)$ e $C = C(H,r)$ tais que
\[
  \lim_{n\to\infty} \Pr\!\bigl[\Gnp \ramsey{r} H\bigr] =
  \begin{cases}
    0 & \text{se } p \le c\, n^{-1/\mdens(H)},\\
    1 & \text{se } p \ge C\, n^{-1/\mdens(H)}.
  \end{cases}
\]
Em particular, $\hat{p}(n) = n^{-1/\mdens(H)}$ é o limiar semi-agudo para
a propriedade $r$-Ramsey de $H$ em $\Gnp$.
\end{theorem}

Uma prova curta e elegante do Teorema~\ref{thm:rodl-rucinski} foi dada por
Nenadov e Steger \cite{nenadov2016short}, usando a técnica de
\emph{containers} de hipergrafos.

% ============================================================
\section{Propriedade Anti-Ramsey em Grafos Aleatórios} \label{sec:antiramsey}
% ============================================================

\begin{definition}[Propriedade anti-Ramsey]
\label{def:antiramsey}
Dados grafos $G$ e $H$, escrevemos $G \rb H$ para denotar a
\textbf{propriedade anti-Ramsey} de $G$ com respeito a $H$:
para toda coloração própria das arestas de $G$
(com número arbitrário de cores),
existe pelo menos uma cópia \emph{arco-íris} de $H$ em $G$,
isto é, uma cópia de $H$ em que cada aresta recebe uma cor distinta.
\end{definition}

\begin{remark}
A propriedade $G \rb H$ é monótona crescente em $G$.
Logo, para qualquer grafo fixo $H$, o grafo aleatório $\Gnp$ admite
um limiar $p_H^{\mathrm{rb}} = p_H^{\mathrm{rb}}(n)$ para essa propriedade
(cf.~Bollobás--Thomason).
\end{remark}

A questão central é determinar $p_H^{\mathrm{rb}}$ em termos de parâmetros
combinatórios de $H$.

\subsection{1-Enunciado: limite superior para o limiar}

O primeiro resultado geral foi obtido por Rödl e Tuza~\cite{rodl1992rainbow}
para ciclos, e depois estendido por Kohayakawa, Konstadinidis e Mota:

\begin{theorem}[Kohayakawa--Konstadinidis--Mota \cite{kohayakawa2014anti}]
\label{thm:1statement}
Seja $H$ um grafo qualquer. Existe uma constante $C > 0$ tal que,
se $p = p(n) \ge C\, n^{-1/\mdens(H)}$, então \whp\
$\Gnp \rb H$.
Em particular, $p_H^{\mathrm{rb}} \le n^{-1/\mdens(H)}$.
\end{theorem}

Em outras palavras, o limiar anti-Ramsey nunca está acima do limiar Ramsey clássico.

\subsection{0-Enunciado: limite inferior — o \emph{framework} de Nenadov et al.}

Para o limite inferior, Nenadov, Person, Škorić e Steger desenvolveram um
arcabouço geral que reduz o problema probabilístico a uma questão determinística:

\begin{theorem}[Nenadov--Person--Škorić--Steger \cite{nenadov2017algorithmic}]
\label{thm:framework}
Seja $H$ um grafo estritamente $2$-balanceado.
Se todo grafo $G$ tal que $G \rb H$ satisfaz $\maxdens(G) > \mdens(H)$,
então existe uma constante $c > 0$ tal que, se $p \le c\, n^{-1/\mdens(H)}$,
então \whp\ $\Gnp \not\rb H$.

Em outras palavras, verificar a condição determinística acima
é suficiente para obter o $0$-enunciado semi-agudo.
\end{theorem}

\begin{definition}[Condição determinística anti-Ramsey]
\label{def:cond-det}
Dado um grafo $H$, dizemos que $H$ \textbf{satisfaz a condição}~(D)
se todo grafo $G$ tal que $G \rb H$ satisfaz $\maxdens(G) > \mdens(H)$.
\end{definition}

\subsection{Resultados para classes específicas de grafos}

\subsubsection*{Ciclos}

\begin{theorem}[Barros--Cavalar--Mota--Parczyk \cite{AntiCycles}]
\label{thm:anticycles}
O limiar $p_{C_\ell}^{\mathrm{rb}}$ para a propriedade anti-Ramsey do ciclo $C_\ell$
satisfaz:
\begin{itemize}
  \item Se $\ell \ge 5$, então
    $p_{C_\ell}^{\mathrm{rb}} = n^{-1/\mdens(C_\ell)} = n^{-(\ell-2)/(\ell-1)}$
    (limiar semi-agudo).
  \item Para $\ell = 4$:
    $p_{C_4}^{\mathrm{rb}} = n^{-3/4} \ll n^{-1/\mdens(C_4)} = n^{-2/3}$.
  \item Para $\ell = 3$: a condição de ser grafo arco-íris é satisfeita
    trivialmente, logo
    $p_{C_3}^{\mathrm{rb}} = n^{-1}$ (limiar para a presença de um triângulo).
\end{itemize}
\end{theorem}

\subsubsection*{Grafos completos}

\begin{theorem}[Kohayakawa--Mota--Parczyk--Schnitzer \cite{AntiComplete}]
\label{thm:anticomplete}
Para o clique $K_k$, o limiar $p_{K_k}^{\mathrm{rb}}$ satisfaz:
\begin{itemize}
  \item Se $k \ge 5$, então
    $p_{K_k}^{\mathrm{rb}} = n^{-1/\mdens(K_k)} = n^{-2/(k+1)}$
    (limiar semi-agudo).
  \item Para $k = 4$:
    $p_{K_4}^{\mathrm{rb}} = n^{-7/15} \ll n^{-1/\mdens(K_4)} = n^{-2/5}$.
  \item Para $k = 3$: trivialmente $p_{K_3}^{\mathrm{rb}} = n^{-1}$.
\end{itemize}
\end{theorem}

\subsubsection*{Caso geral: 2-densidade alta}

\begin{theorem}[Kuperwasser \cite{kuperwasser2026anti}]
\label{thm:kuperwasser}
Seja $H$ um grafo com $\mdens(H) \ge 19$.
Existem constantes $c_0 < c_1$ tais que
\[
  \lim_{n\to\infty}\Pr\!\bigl[\Gnp \rb H\bigr] =
  \begin{cases}
    0 & \text{se } p \le c_0\, n^{-1/\mdens(H)},\\
    1 & \text{se } p \ge c_1\, n^{-1/\mdens(H)}.
  \end{cases}
\]
\end{theorem}

O Teorema~\ref{thm:kuperwasser} é provado verificando a condição determinística
do Teorema~\ref{thm:framework} através do seguinte resultado:

\begin{theorem}[Kuperwasser \cite{kuperwasser2026anti}]
\label{thm:kuperwasser-det}
Todo grafo $G$ com $\maxdens(G) \ge 18$ admite uma coloração própria das arestas
tal que todo subgrafo arco-íris tem $2$-densidade estritamente menor que $\maxdens(G)$.
Em particular, se $\mdens(H) > 18$, então $H$ satisfaz a condição~\emph{(D)}.
\end{theorem}

\subsection{A tabela de limiares conhecidos}

A tabela a seguir resume o estado atual dos limiares para a propriedade anti-Ramsey,
adaptada de \cite{behague2025thresholds}:

\begin{center}
\renewcommand{\arraystretch}{1.4}
\begin{tabular}{lll}
\hline
\textbf{Grafo $H$} & \textbf{Limiar $p_H^{\mathrm{rb}}$} & \textbf{Referência}\\
\hline
$K_\ell$, $\ell \ge 5$ & $n^{-2/(\ell+1)}$ (semi-agudo) & \cite{AntiComplete}\\
$K_4$                  & $n^{-7/15} \ll n^{-1/m_2(K_4)}$& \cite{AntiComplete}\\
$K_3$                  & $n^{-1}$                        & trivial\\
$C_\ell$, $\ell \ge 5$ & $n^{-(\ell-2)/(\ell-1)}$ (semi-agudo) & \cite{AntiCycles}\\
$C_4$                  & $n^{-3/4} \ll n^{-1/m_2(C_4)}$ & \cite{AntiCycles}\\
$C_3$                  & $n^{-1}$                        & trivial\\
$H$ com $m_2(H)\ge 19$ & $n^{-1/m_2(H)}$ (semi-agudo)   & \cite{kuperwasser2026anti}\\
\hline
\end{tabular}
\end{center}

\subsection{Questão em aberto: grafos de cintura grande}

Os resultados acima revelam um padrão: o limiar anti-Ramsey diverge de
$n^{-1/\mdens(H)}$ exatamente para grafos ``pequenos'' de cintura baixa
($K_3$, $K_4$, $C_4$). Essa observação motiva a seguinte questão central
deste projeto de doutorado:

\begin{question}
\label{q:cintura}
Seja $H$ um grafo com cintura (comprimento do menor ciclo) estritamente maior que~$4$.
É verdade que $p_H^{\mathrm{rb}} = n^{-1/\mdens(H)}$?
Equivalentemente, $H$ sempre satisfaz a condição~\emph{(D)}
(Definição~\ref{def:cond-det}) quando $\mathrm{girth}(H) > 4$?
\end{question}

% ============================================================
\section{Propriedade Ramsey Restrita (Constrained Ramsey)} \label{sec:cram}
% ============================================================

Este problema, estudado em \cite{behague2025thresholds}, generaliza tanto
a propriedade de Ramsey clássica quanto a anti-Ramsey.

\begin{definition}[Propriedade Ramsey restrita]
\label{def:cram}
Dados grafos $H_1$ e $H_2$, escrevemos $G \cram (H_1, H_2)$ para denotar
a \textbf{propriedade Ramsey restrita} (\emph{constrained Ramsey}):
para toda coloração própria das arestas de $G$,
existe uma cópia monocromática de $H_1$ \emph{ou} uma cópia arco-íris de $H_2$.
\end{definition}

\begin{remark}
Quando $H_2$ é uma floresta, qualquer coloração própria tem trivialmente uma
cópia arco-íris de $H_2$ (desde que o grafo seja grande o suficiente). Logo,
o caso interessante é quando $H_2$ não é uma floresta.

Quando $H_1 = $ estrela com $k$ arestas $K_{1,k}$ e $H_2 = H$, a propriedade
$G \cram (K_{1,k}, H)$ equivale a dizer que toda coloração \emph{localmente
$(k-1)$-limitada} de $G$ (em que cada cor aparece no máximo $k-1$ vezes por vértice)
contém uma cópia arco-íris de $H$. Assim,
a propriedade anti-Ramsey é o caso particular $k = 2$
(colorações próprias = colorações localmente $1$-limitadas).
\end{remark}

\subsection{Resultados principais de Behague et al.}

\begin{theorem}[Behague--Hancock--Hyde--Letzter--Morrison \cite{behague2025thresholds}]
\label{thm:behague-1-stmt}
Sejam $k \ge 2$ e $H$ um grafo qualquer.
Existe uma constante $C > 0$ tal que, se $p \ge C\, n^{-1/\mdens(H)}$, então
\whp\ $\Gnp \cram (K_{1,k}, H)$.
\end{theorem}

O Teorema~\ref{thm:behague-1-stmt} generaliza o Teorema~\ref{thm:1statement}.

\begin{definition}[Densidades relevantes para Ramsey restrito]
Para o problema restrito com $H_1 = K_{1,k}$, definem-se:
\begin{align*}
  m_a(H) &:= \inf\!\bigl\{\maxdens(G) : G \cram (K_{1,k}, H)\bigr\}\\
         &\phantom{:}= \text{densidade máxima do menor grafo} \ G\ \text{que força }\\
         &\phantom{= \text{densidade máxima do menor grafo} G \text{que força}} \text{Ramsey restrito para } (K_{1,k}, H).
\end{align*}
O limiar candidato natural é então $n^{-1/m_a(H)}$.
\end{definition}

\begin{theorem}[Behague et al.\ \cite{behague2025thresholds} -- $0$-enunciado geral]
\label{thm:behague-0-stmt}
Sejam $k \ge 3$ e $H$ um grafo que não é floresta.
\begin{enumerate}
  \item[(a)] Se $k \ge 4$, \emph{ou} $\mdens(H) \ne 2$, \emph{ou}
    $\mdens(H) = 2$ e $H$ contém um subgrafo estritamente $2$-balanceado
    $J \ne K_3$ com $\mdens(J) = 2$,
    então existem $c, C > 0$ tais que
    \[
      \lim_{n\to\infty} \Pr\!\bigl[\Gnp \cram (K_{1,k}, H)\bigr] =
      \begin{cases}
        0 & \text{se } p \le c\, n^{-1/\mdens(H)},\\
        1 & \text{se } p \ge C\, n^{-1/\mdens(H)}.
      \end{cases}
    \]
  \item[(b)] Se $k = 3$ e $\mdens(H) = 2$ e o único subgrafo estritamente
    $2$-balanceado de $H$ com $2$-densidade~$2$ é $K_3$ (e.g.\ $H = K_3$),
    então o limiar é mais fraco: coarse threshold em $n^{-1/2}$.
\end{enumerate}
\end{theorem}

% ============================================================
\section{Propriedade Ramsey Canônica Assimétrica} \label{sec:canonico}
% ============================================================

A teoria canônica de Ramsey estuda estruturas que emergem em colorações
\emph{arbitrárias} (sem restrição de número de cores).

\begin{definition}[Coloração canônica]
\label{def:canonica}
Seja $c : E(G) \to \mathbb{N}$ uma coloração arbitrária das arestas de $G$,
com uma ordenação total $v_1 < v_2 < \cdots < v_n$ dos vértices de $G$.
Um subgrafo $H \subseteq G$ possui uma \textbf{coloração canônica} se
satisfaz um dos seguintes padrões:
\begin{enumerate}
  \item \textbf{Monocromático:} todas as arestas de $H$ têm a mesma cor.
  \item \textbf{Arco-íris:} todas as arestas de $H$ têm cores distintas.
  \item \textbf{Lexicográfico (min-lex):} a cor de $v_iv_j$ (com $v_i < v_j$)
    depende unicamente de $v_i$.
  \item \textbf{Lexicográfico (max-lex):} a cor de $v_iv_j$ (com $v_i < v_j$)
    depende unicamente de $v_j$.
\end{enumerate}
\end{definition}

\begin{theorem}[Erdős--Rado \cite{erdos1950combinatorial}]
Para qualquer inteiro $\ell \ge 1$, existe $N = N(\ell)$ tal que toda coloração
arbitrária das arestas de $K_N$ contém uma cópia canônica de $K_\ell$.
\end{theorem}

O resultado de Erdős--Rado foi trazido para o contexto aleatório por Kamčev e Schacht:

\begin{theorem}[Kamčev--Schacht \cite{kamcev2023canonical}]
\label{thm:kamcev}
Para todo $\ell \ge 4$, o limiar para a propriedade de que toda coloração
arbitrária de $\Gnp$ contenha uma cópia canônica de $K_\ell$ é
$\hat{p}(n) = n^{-2/(\ell+1)} = n^{-1/\mdens(K_\ell)}$.
\end{theorem}

\subsection{Cenário assimétrico}

\begin{definition}[Propriedade Ramsey canônica assimétrica]
\label{def:assimetrico}
Dados grafos $H_1$, $H_2$ e $H_3$, escrevemos
\[
  G \xrightarrow{\mathrm{can}} (H_1, H_2, H_3)
\]
para denotar a seguinte propriedade: para toda coloração arbitrária das
arestas de $G$, existe pelo menos
\begin{itemize}
  \item uma cópia \textbf{monocromática} de $H_1$, ou
  \item uma cópia \textbf{arco-íris} de $H_2$, ou
  \item uma cópia \textbf{lexicográfica} de $H_3$.
\end{itemize}
\end{definition}

\begin{question}
\label{q:assimetrico}
Qual é o comportamento da função limiar $\hat{p}(n)$ para a propriedade
$\Gnp \xrightarrow{\mathrm{can}} (H_1, H_2, H_3)$
quando $H_1$, $H_2$ e $H_3$ não são todos iguais?
Em particular, como o limiar depende das densidades $\mdens(H_i)$?
\end{question}

% ============================================================
\section{O Caso Assimétrico \texorpdfstring{$(K_3, K_4, K_5)$}{(K3, K4, K5)}}
% ============================================================

Queremos determinar o limiar exato para a propriedade $\Gnp \xrightarrow{\mathrm{can}} (K_3,K_4,K_5)$.
Note que, para $p \ll n^{-1/\mdens(K_3)} = n^{-1/2}$, o grafo aleatório \whp\ não possui triângulos, tornando a propriedade trivialmente falsa (através de uma coloração arco-íris, por exemplo).
Por outro lado, para $p \gg n^{-1/\mdens(K_5)} = n^{-1/3}$, a existência de um $K_5$ canônico é garantida \whp, o que força a existência de uma das estruturas proibidas.

O resultado principal que buscamos formalizar é que o limiar coincide com o do clique mais denso.
Para a prova do $0$-enunciado, precisaremos das seguintes definições estruturais:

\begin{definition}[$H$-fechamento]
\label{def:H-fechado}
Dado um grafo $G$ e um grafo $H$, dizemos que:
\begin{itemize}
  \item Uma aresta $e \in E(G)$ é \textbf{$H$-fechada} se pertence a pelo menos duas cópias de $H$ em $G$.
  \item Uma cópia $H' \subseteq G$ de $H$ é \textbf{$H$-fechada} se pelo menos três de suas arestas são $H$-fechadas.
  \item O grafo $G$ é \textbf{$H$-fechado} se todo vértice e toda aresta de $G$ pertencem a pelo menos uma cópia de $H$, e toda cópia de $H$ em $G$ é $H$-fechada.
\end{itemize}
\end{definition}

\begin{definition}[$H$-Bloco]
\label{def:H-bloco}
Um grafo $G$ é um \textbf{$H$-bloco} se $G$ é $H$-fechado e, para todo subconjunto próprio não-vazio de arestas $E' \subsetneq E(G)$, existe uma cópia $H'$ de $H$ em $G$ tal que $E(H') \cap E' \ne \emptyset$ e $E(H') \setminus E' \ne \emptyset$.
\end{definition}

\begin{remark}
A chave do framework de Nenadov et al.\ é que, quando $p \ll n^{-1/\maxdens(H)}$, os $H$-blocos presentes em $\Gnp$ têm tamanho $O(1)$ \whp.
\end{remark}

\begin{theorem}
  Existe um limiar semi-agudo para a propriedade $\Gnp \xrightarrow{\mathrm{can}} (K_3,K_4,K_5)$ em $\hat{p}(n) = n^{-1/3}$. Ou seja, existem constantes $c, C > 0$ tais que:
  \[
  \lim_{n \to \infty} \mathbb{P}(\Gnp \xrightarrow{\mathrm{can}} (K_3,K_4,K_5)) = 
  \begin{cases}
    1 & \text{se } p \ge C n^{-1/3}, \\
    0 & \text{se } p \le c n^{-1/3}.  
  \end{cases}
  \]
\end{theorem}

\begin{proof}[Esbo\c{c}o da Prova]
  O $1$-enunciado segue do Teorema de Kam\v{c}ev--Schacht~\cite{kamcev2023canonical}.

  Para o $0$-enunciado, descrevemos o algoritmo \textsc{Colora\c{c}\~{a}oCan} que, dado $G$ com $\maxdens(G) \le 3$, constr\'{o}i explicitamente uma colora\c{c}\~{a}o das arestas de $G$ sem $K_3$ monocrom\'{a}tico, $K_4$ arco-\'{i}ris ou $K_5$ lexicogr\'{a}fico.

  \textbf{Passo 1 --- Colora\c{c}\~{a}o Lexicogr\'{a}fica Inicial.}
  Fixamos uma ordena\c{c}\~{a}o total $\sigma: V(G) \to [n]$ e uma atribui\c{c}\~{a}o $\varphi: V(G) \to \mathbb{N}$.
  Definimos a colora\c{c}\~{a}o min-lex $\chi: E(G) \to \mathbb{N}$ por
  \[ \chi(uv) \;:=\; \varphi\!\bigl(\min_\sigma\{u,v\}\bigr), \]
  ou seja, a cor de cada aresta \'{e} determinada pelo v\'{e}rtice de menor \'{i}ndice segundo $\sigma$.
  Por constru\c{c}\~{a}o, $\chi$ \'{e} uma colora\c{c}\~{a}o min-lex e satisfaz:
  \begin{itemize}
    \item n\~{a}o existe $K_3$ monocrom\'{a}tico sob $c$;
    \item n\~{a}o existe $K_4$ arco-\'{i}ris sob $c$.
  \end{itemize}
  O problema \'{e} que toda c\'{o}pia de $K_5$ em $G$ \'{e} lexicogr\'{a}fica (min-lex) sob $\chi$ --- exatamente a estrutura proibida que precisamos eliminar.

  \textbf{Passo 2 --- Remo\c{c}\~{a}o de Arestas Livres.}
  Removemos de $G$ todas as arestas que n\~{a}o pertencem a nenhuma c\'{o}pia de $K_5$.
  Seja $G' \subseteq G$ o grafo resultante. As arestas removidas mant\^{e}m a colora\c{c}\~{a}o $\chi$ do Passo~1 e n\~{a}o causam problemas.

  \textbf{Passo 3 --- Decomposi\c{c}\~{a}o em $K_5$-Blocos.}
  Decompomos $G'$ em seus $K_5$-blocos (Defini\c{c}\~{a}o~\ref{def:H-bloco}).
  Como $\maxdens(G) \le \mdens(K_5) = 3$, cada $K_5$-bloco tem tamanho $O(1)$.

  \textbf{Passo 4 --- Recolora\c{c}\~{a}o dos Blocos.}
  Para cada $K_5$-bloco $B$, substitu\'{i}mos $\chi|_{E(B)}$ por uma nova colora\c{c}\~{a}o $\chi'$ (constru\'{i}da indutivamente) que evita $K_5$ lexicogr\'{a}fico, sem introduzir $K_3$ monocrom\'{a}tico ou $K_4$ arco-\'{i}ris.
  A exist\^{e}ncia de tal colora\c{c}\~{a}o em cada bloco \'{e} garantida pelo Lema~\ref{lem:0-statement-K5} abaixo.
\end{proof}

\begin{lemma}
  \label{lem:0-statement-K5}
  Seja $B$ um $K_5$-bloco com $\maxdens(B) \le 3$. Ent\~{a}o existe uma colora\c{c}\~{a}o $c: E(B) \to \mathbb{N}$ tal que:
  \begin{enumerate}
    \item[(i)] $c$ n\~{a}o possui $K_3$ monocrom\'{a}tico;
    \item[(ii)] $c$ n\~{a}o possui $K_4$ arco-\'{i}ris;
    \item[(iii)] $c$ n\~{a}o possui $K_5$ lexicogr\'{a}fico.
  \end{enumerate}
\end{lemma}

\begin{proof}

\noindent\textit{A colora\c{c}\~{a}o auxiliar de $K_5$.}

Seja $B$ um $K_5$-bloco com $\maxdens(B) \le 3$.
Consideramos uma ordena\c{c}\~{a}o $\sigma$ dos v\'{e}rtices de $B$ e a colora\c{c}\~{a}o lexicogr\'{a}fica min-lex
$\chi(uv) := \varphi(\min_\sigma\{u,v\})$.

O ponto de partida \'{e} uma \'{u}nica c\'{o}pia $H_0 \cong K_5$ com v\'{e}rtices ordenados $v_1 < v_2 < v_3 < v_4 < v_5$.
Sob a colora\c{c}\~{a}o min-lex, as cores das arestas s\~{a}o distribu\'{i}das como segue:

\medskip
\medskip
\begin{center}
\begin{tikzpicture}[scale=1.3,
  vx/.style={circle,draw,fill=white,minimum size=16pt,font=\small}]
  % K5 vertices as regular pentagon
  \node[vx] (1) at (90:1.5)  {$v_1$};
  \node[vx] (2) at (18:1.5)  {$v_2$};
  \node[vx] (3) at (-54:1.5) {$v_3$};
  \node[vx] (4) at (-126:1.5){$v_4$};
  \node[vx] (5) at (162:1.5) {$v_5$};
  % color a = phi(v1): edges 12,13,14,15
  \draw[line width=1.8pt, red]    (1)--(2);
  \draw[line width=1.8pt, red]    (1)--(3);
  \draw[line width=1.8pt, red]    (1)--(4);
  \draw[line width=1.8pt, red]    (1)--(5);
  % color b = phi(v2): edges 23,24,25
  \draw[line width=1.8pt, blue]   (2)--(3);
  \draw[line width=1.8pt, blue]   (2)--(4);
  \draw[line width=1.8pt, blue]   (2)--(5);
  % color c = phi(v3): edges 34,35
  \draw[line width=1.8pt, teal]   (3)--(4);
  \draw[line width=1.8pt, teal]   (3)--(5);
  % color d = phi(v4): edge 45
  \draw[line width=1.8pt, orange] (4)--(5);
\end{tikzpicture}
\end{center}

\noindent
\'{E} f\'{a}cil verificar que toda c\'{o}pia de $K_5$ sob $\chi$ \'{e} lexicogr\'{a}fica (min-lex).

Definimos agora uma \textbf{colora\c{c}\~{a}o auxiliar de duas cores} $\chi': E(B) \to \{\text{claro}, \text{escuro}\}$
(representadas como vermelho e azul nos diagramas) que servir\'{a} de base para a recolora\c{c}\~{a}o dos $K_5$-blocos.
Sobre o $K_5$ inicial $H_0$, $\chi'$ \'{e} definida identificando os $10$ arestas de $K_5$ com as $10$ arestas de
duas c\'{o}pias de $C_5$: o \textbf{pent\'{a}gono externo} (as arestas $v_1v_2, v_2v_3, v_3v_4, v_4v_5, v_5v_1$)
e o \textbf{pent\'{a}gono interno} (as diagonais $v_1v_3, v_3v_5, v_5v_2, v_2v_4, v_4v_1$).

\begin{center}
\begin{tikzpicture}[scale=1.3,
  vx/.style={circle,draw,fill=white,minimum size=16pt,font=\small}]
  \node[vx] (1) at (90:1.5)  {$v_1$};
  \node[vx] (2) at (18:1.5)  {$v_2$};
  \node[vx] (3) at (-54:1.5) {$v_3$};
  \node[vx] (4) at (-126:1.5){$v_4$};
  \node[vx] (5) at (162:1.5) {$v_5$};
  % outer C5 (pentagon) -- claro = vermelho
  \draw[line width=2pt, red]  (1)--(2)--(3)--(4)--(5)--(1);
  % inner C5 (pentagram) -- escuro = azul
  \draw[line width=2pt, blue] (1)--(3)--(5)--(2)--(4)--(1);
\end{tikzpicture}
\end{center}

\noindent
Verifica-se diretamente que $\chi|_{H_0}$ \'{e} uma colora\c{c}\~{a}o pr\'{o}pria de $H_0$ (todo v\'{e}rtice tem grau $2$ em cada $C_5$),
n\~{a}o possui $K_3$ monocrom\'{a}tico, n\~{a}o possui $K_4$ arco-\'{i}ris, e n\~{a}o possui $K_5$ lexicogr\'{a}fico.


\medskip
\noindent\textit{Fun\c{c}\~{a}o de penalidade e tipos de expans\~{a}o.}

Para controlar a estrutura dos $K_5$-blocos durante a constru\c{c}\~{a}o indutiva,
introduzimos a seguinte medida adaptada de~\cite{AntiComplete}:

\begin{definition}[Fun\c{c}\~{a}o de penalidade]
Dado um grafo $G$, definimos
\[
  p(G) \;:=\; e(G) - 3\,v(G) + 5.
\]
\end{definition}

A fun\c{c}\~{a}o $p$ \'{e} normalizada de forma que o $K_5$ inicial tenha penalidade nula:
\[
  p(K_5) = 10 - 3 \cdot 5 + 5 = 0.
\]
Al\'{e}m disso, a condi\c{c}\~{a}o de densidade $\maxdens(G) \le 3$ garante que $e(G) \le 3\,v(G)$, o que implica
$p(G) \le 5$ para qualquer subgrafo de $B$. Essa limita\c{c}\~{a}o superior \'{e} crucial para controlar a 
complexidade estrutural do bloco.

A constru\c{c}\~{a}o de um $K_5$-bloco $B$ pode ser vista como um processo indutivo. Come\c{c}amos com uma 
c\'{o}pia base $H_0 \cong K_5$ e, a cada passo, anexamos novos v\'{e}rtices e arestas para formar novas 
c\'{o}pias de $K_5$. Seguindo as estrat\'{e}gias desenvolvidas em~\cite{nenadov2017algorithmic} e \cite{AntiComplete},
classificamos essas extens\~{o}es em quatro tipos estruturais (ver analogia aos casos de expans\~{a}o):

\begin{itemize}
  \item \textbf{E0 (Expans\~{a}o regular):} Adicionam-se $3$ novos v\'{e}rtices e $9$ novas arestas que, juntamente 
    com uma aresta j\'{a} existente, formam um novo $K_5$. 
    Neste caso, a varia\c{c}\~{a}o da penalidade \'{e} $\Delta p = 9 - 3(3) = 0$, ou seja, a penalidade do 
    grafo se mant\'{e}m inalterada.

  \item \textbf{E1 (Expans\~{a}o sobre tri\^{a}ngulo):} Adicionam-se $2$ novos v\'{e}rtices e $7$ novas arestas que, 
    apoiadas em um tri\^{a}ngulo preexistente, completam um novo $K_5$. 
    Aqui, o incremento na penalidade \'{e} $\Delta p = 7 - 3(2) = 1$.

  \item \textbf{E2 (Expans\~{a}o sobre $K_4$):} Adiciona-se apenas $1$ novo v\'{e}rtice ligado a um $K_4$ 
    existente atrav\'{e}s de $4$ novas arestas, completando um $K_5$. 
    O incremento \'{e} $\Delta p = 4 - 3(1) = 1$.

  \item \textbf{U1 (Uni\~{a}o sobre $K_5^-$):} N\~{a}o se adicionam v\'{e}rtices; apenas $1$ nova aresta \'{e} 
    inserida conectando dois v\'{e}rtices de um $K_5$ quase completo ($K_5 \setminus \{e\}$) j\'{a} 
    presente na estrutura. 
    A varia\c{c}\~{a}o \'{e} $\Delta p = 1 - 3(0) = 1$.
\end{itemize}

Uma observa\c{c}\~{a}o fundamental \'{e} que a expans\~{a}o \textbf{E0} (a \'{u}nica que \'{e} considerada "regular") 
preserva o valor de $p(G)$. Por outro lado, as expans\~{o}es degeneradas (\textbf{E1}, \textbf{E2} e \textbf{U1}) 
aumentam a penalidade estritamente em $1$. Como iniciamos com $p(H_0) = 0$ e sabemos que $p(B) \le 5$, conclu\'{i}mos 
que o n\'{u}mero total de ocorr\^{e}ncias dos casos E1, E2 e U1 na constru\c{c}\~{a}o do bloco \'{e} uniformemente 
limitado (ocorrem no m\'{a}ximo $5$ vezes em todo o processo).

\medskip
\noindent\textit{Extens\~{a}o da colora\c{c}\~{a}o auxiliar para cada tipo de expans\~{a}o.}
\medskip

\noindent\textbf{Caso E0 (expans\~{a}o regular).}
Em E0, uma nova c\'{o}pia de $K_5$ \'{e} colada \`{a} estrutura existente por uma aresta $e$ j\'{a} colorida.
Identificamos os dois $C_5$ na nova c\'{o}pia de modo que $e$ perten\c{c}a ao $C_5$ da sua cor.
A colorac\~{a}o \'{e} compat\'{i}vel com o bloco anterior.
No diagrama abaixo, dois $K_5$ s\~{a}o colados pela aresta $1$--$2$ (vermelha):

\medskip
\begin{center}
\begin{tikzpicture}[scale=1.0,
  vx/.style={circle,draw,fill=white,minimum size=13pt,font=\tiny\strut}]
  \node[vx] (v1) at (0, 0.9)   {$1$};
  \node[vx] (v2) at (0,-0.9)   {$2$};
  \node[vx] (v5) at (-1.4, 1.6){$5$};
  \node[vx] (v4) at (-2.0, 0)  {$4$};
  \node[vx] (v3) at (-1.4,-1.6){$3$};
  \node[vx] (v8) at ( 1.4, 1.6){$8$};
  \node[vx] (v7) at ( 2.0, 0)  {$7$};
  \node[vx] (v6) at ( 1.4,-1.6){$6$};
  % Left K5: Red outer C5: 1-5-4-3-2-1
  \draw[red, line width=1.8pt]  (v1)--(v5)--(v4)--(v3)--(v2);
  % Left K5: Blue inner C5: 1-4-2-5-3-1
  \draw[blue, line width=1.8pt] (v1)--(v4) (v4)--(v2) (v2)--(v5) (v5)--(v3) (v3)--(v1);
  % Right K5: Red outer C5: 1-8-7-6-2-1
  \draw[red, line width=1.8pt]  (v1)--(v8)--(v7)--(v6)--(v2);
  % Right K5: Blue inner C5: 1-7-2-8-6-1
  \draw[blue, line width=1.8pt] (v1)--(v7) (v7)--(v2) (v2)--(v8) (v8)--(v6) (v6)--(v1);
  % Shared edge highlighted
  \draw[red, line width=3pt] (v1)--(v2);
  \node[font=\tiny] at (0.25,0) {$e$};
\end{tikzpicture}
\end{center}

\medskip
\noindent\textbf{Caso E1 (expans\~{a}o sobre tri\^{a}ngulo).}
Em E1, dois novos v\'{e}rtices $u, v$ com a aresta $uv$ s\~{a}o ligados a um tri\^{a}ngulo $\{a,b,c\}$ existente.
Identificamos um $P_2$ monocrom\'{a}tico no tri\^{a}ngulo (sempre existe por paridade) e extendemos
ao $C_5$ da sua cor; as cinco arestas restantes formam o $C_5$ da outra cor.
Se o $P_2$ vermelho \'{e} $a$--$b$--$c$, ent\~{a}o: $C_5$ vermelho $= a\text{-}b\text{-}c\text{-}u\text{-}v\text{-}a$;
$C_5$ azul $= a\text{-}c\text{-}v\text{-}b\text{-}u\text{-}a$.

\medskip
\begin{center}
\begin{tikzpicture}[scale=1,
  vx/.style={circle,draw,fill=white,minimum size=13pt,font=\tiny\strut},
  ov/.style={circle,draw,fill=gray!25,minimum size=13pt,font=\tiny\strut}]

%================================================
% K5 com coordenadas explícitas
%================================================

\node[ov] (A) at (0,1.8)    {$a$};
\node[ov] (B) at (1.7,0.6)  {$b$};
\node[ov] (C) at (1.0,-1.4) {$c$};
\node[vx] (U) at (-1.0,-1.4) {$u$};
\node[vx] (V) at (-1.7,0.6) {$v$};

% vermelho: a-b-c-u-v-a
\draw[red, line width=1.8pt]
(A)--(B)--(C)--(U)--(V)--(A);

% azul: a-c-v-b-u-a
\draw[blue, line width=1.8pt]
(A)--(C)
(C)--(V)
(V)--(B)
(B)--(U)
(U)--(A);

%================================================
% vértices extras x,y à direita
%================================================

\node[vx] (X) at (4.0,0.8)  {$x$};
\node[vx] (Y) at (4.0,-0.8) {$y$};

% aresta xy
\draw[black, line width=1.5pt] (X)--(Y);

% ligações ao triângulo abc
\draw[dashed,gray]
(X)--(A)
(X)--(B)
(X)--(C);

\draw[dashed,gray]
(Y)--(A)
(Y)--(B)
(Y)--(C);

\end{tikzpicture}
\end{center}
\medskip
\noindent\textbf{Caso E2 (expans\~{a}o sobre $K_4$).}
Em E2, um novo v\'{e}rtice $v$ \'{e} ligado a todos os v\'{e}rtices de um $K_4$ existente $\{a,b,c,d\}$.
O $K_4$ possui dois $P_3$ monocrom\'{a}ticos; tomamos o $P_3$ vermelho $a$--$b$--$c$--$d$ e extendemos:
$C_5$ vermelho $= a\text{-}b\text{-}c\text{-}d\text{-}v\text{-}a$;
$C_5$ azul $= a\text{-}c\text{-}v\text{-}b\text{-}d\text{-}a$.

\medskip
\begin{center}
\begin{tikzpicture}[scale=1,
  vx/.style={circle,draw,fill=white,minimum size=13pt,font=\tiny\strut},
  ov/.style={circle,draw,fill=gray!25,minimum size=13pt,font=\tiny\strut}]

%================================================
% K5
%================================================

\node[ov] (A) at (0,1.8)     {$a$};
\node[ov] (B) at (1.7,0.6)   {$b$};
\node[ov] (C) at (1.0,-1.4)  {$c$};
\node[vx] (U) at (-1.0,-1.4) {$u$};
\node[vx] (V) at (-1.7,0.6)  {$v$};

% vermelho: a-b-c-u-v-a
\draw[red, line width=1.8pt]
(A)--(B)--(C)--(U)--(V)--(A);

% azul: a-c-v-b-u-a
\draw[blue, line width=1.8pt]
(A)--(C)
(C)--(V)
(V)--(B)
(B)--(U)
(U)--(A);

%================================================
% vértice extra à direita
%================================================

\node[vx] (X) at (4.2,0) {$x$};

% x ligado ao K4 abcv
\draw[dashed,gray]
(X)--(A)
(X)--(B)
(X)--(C)
(X)--(V);

\end{tikzpicture}
\end{center}
\medskip

\begin{proposition}[Propriedade de equil\'{i}brio de cores]
\label{prop:equilibrio}
Em cada etapa das expans\~{o}es $E_0$, $E_1$ e $E_2$, \'{e} poss\'{i}vel construir uma colora\c{c}\~{a}o
dos $K_5$-blocos tal que, para toda c\'{o}pia de $K_5$ no bloco e todo v\'{e}rtice $w$ dessa c\'{o}pia,
$w$ possui pelo menos um vizinho de cor \textbf{claro} e pelo menos um vizinho de cor \textbf{escuro}
dentro da c\'{o}pia de $K_5$.
\end{proposition}

\noindent
Antes de provar o Caso U1, mostramos que a Proposi\c{c}\~{a}o~\ref{prop:equilibrio} \'{e} suficiente
para concluir a prova do Lema~\ref{lem:0-statement-K5}.

\begin{remark}[A colora\c{c}\~{a}o $\chi''$ e as tr\^{e}s condi\c{c}\~{o}es do Lema]
\label{rem:equilibrio-suficiente}
Dadas $\chi$ e $\chi'$ como acima, definimos a \textbf{colora\c{c}\~{a}o final}
\[
  \chi'': E(B) \;\longrightarrow\; \mathbb{N} \times \{\text{claro}, \text{escuro}\}
\]
por $\chi''(e) := \bigl(\chi(e),\, \chi'(e)\bigr)$, ou seja, cada aresta $e$ recebe como cor
o par formado pela sua cor lexicogr\'{a}fica (dada por $\chi$) e pela sua cor auxiliar (dada por $\chi'$).

Suponha que $\chi'$ satisfaz a Proposi\c{c}\~{a}o~\ref{prop:equilibrio}, isto \'{e},
todo v\'{e}rtice $w$ de qualquer $K_5$ do bloco possui ao menos um vizinho com
$\chi'(e) = \text{claro}$ e ao menos um vizinho com $\chi'(e) = \text{escuro}$ nessa c\'{o}pia.
Mostraremos que $\chi''$ evita as tr\^{e}s estruturas proibidas.

\begin{itemize}
  \item \textbf{N\~{a}o existe $K_5$ lexicogr\'{a}fico sob $\chi''$.}
  Sob $\chi$, toda c\'{o}pia de $K_5$ \'{e} lexicogr\'{a}fica: fixado o v\'{e}rtice m\'{i}nimo
  $v_1 < v_2 < \cdots < v_5$, todas as quatro arestas $v_1v_j$ ($j=2,3,4,5$) t\^{e}m cor
  $\chi(v_1v_j) = \varphi(v_1)$, a mesma para todas.
  Uma c\'{o}pia de $K_5$ \'{e} min-lex sob $\chi''$ se e somente se todas as arestas
  $v_1v_j$ t\^{e}m a mesma cor $\chi''(v_1v_j) = (\varphi(v_1), \chi'(v_1v_j))$.
  Isso ocorreria somente se todas as quatro arestas de $v_1$ tivessem $\chi' = \text{claro}$
  ou todas $\chi' = \text{escuro}$.
  Pela Proposi\c{c}\~{a}o~\ref{prop:equilibrio}, $v_1$ possui ao menos uma aresta com $\chi' = \text{claro}$
  e ao menos uma com $\chi' = \text{escuro}$. Logo, as arestas de $v_1$ n\~{a}o s\~{a}o todas iguais em
  $\chi''$, e o $K_5$ n\~{a}o \'{e} min-lex. O caso max-lex \'{e} an\'{a}logo (v\'{e}rtice m\'{a}ximo).

  \item \textbf{N\~{a}o existe $K_4$ arco-\'{i}ris sob $\chi''$.}
  Seja $\{w, u_1, u_2, u_3\}$ uma c\'{o}pia de $K_4$ no bloco.
  Sob $\chi$, as tr\^{e}s arestas $wu_1, wu_2, wu_3$ incidentes ao v\'{e}rtice m\'{i}nimo $w$
  t\^{e}m todas a mesma cor $\chi(wu_i) = \varphi(w)$.
  Sob $\chi''$, essas arestas t\^{e}m cores $(\varphi(w), \chi'(wu_i))$ para $i = 1, 2, 3$.
  Pela Proposi\c{c}\~{a}o~\ref{prop:equilibrio} aplicada a qualquer $K_5$ que contenha $\{w, u_1, u_2, u_3\}$,
  $w$ tem ao menos uma aresta clara e uma escura entre $u_1, u_2, u_3$.
  Portanto, duas dessas tr\^{e}s arestas t\^{e}m $\chi'$ iguais, logo t\^{e}m $\chi''$ iguais.
  O $K_4$ n\~{a}o \'{e} arco-\'{i}ris.

  \item \textbf{N\~{a}o existe $K_3$ monocrom\'{a}tico sob $\chi''$.}
  Seja $\{x, y, z\}$ um tri\^{a}ngulo com $x < y < z$ segundo $\sigma$.
  As tr\^{e}s arestas t\^{e}m cores $\chi''(xy) = (\varphi(x), \cdot)$,
  $\chi''(xz) = (\varphi(x), \cdot)$ e $\chi''(yz) = (\varphi(y), \cdot)$.
  Como $x \ne y$ (v\'{e}rtices distintos), $\varphi(x) \ne \varphi(y)$ em geral, logo
  $\chi''(xy) \ne \chi''(yz)$, e o tri\^{a}ngulo n\~{a}o \'{e} monocrom\'{a}tico.
  (Mais precisamente, $\chi$ \'{e} min-lex e $\varphi$ \'{e} injetiva nos r\'{o}tulos dos v\'{e}rtices,
  garantindo que arestas com v\'{e}rtices m\'{i}nimos distintos t\^{e}m cores $\chi$ distintas.)
\end{itemize}

Em conclus\~{a}o, provada a Proposi\c{c}\~{a}o~\ref{prop:equilibrio},
a colora\c{c}\~{a}o $\chi''$ \'{e} v\'{a}lida para o Lema~\ref{lem:0-statement-K5}.
\end{remark}

\medskip
\noindent\textbf{Caso U1 (uni\~{a}o sobre $K_5^-$).}
Em U1, adicionamos uma \'{u}nica aresta $ae$ que completa um $K_5^- = \{a,b,c,d,e\} \setminus \{ae\}$
j\'{a} existente no bloco, formando um novo $K_5$.

As nove arestas do $K_5^-$ j\'{a} est\~{a}o coloridas (claro/escuro) pelas etapas anteriores.
Pela hip\'{o}tese indutiva (Proposi\c{c}\~{a}o~\ref{prop:equilibrio}),
cada v\'{e}rtice de qualquer $K_5$ anterior que contenha v\'{e}rtices de $\{a,b,c,d,e\}$
possui pelo menos um vizinho claro e um vizinho escuro.

Em particular, no $K_5^-$:
\begin{itemize}
  \item $a$ possui tr\^{e}s vizinhos $\{b,c,d\}$ j\'{a} coloridos e, pela Proposi\c{c}\~{a}o~\ref{prop:equilibrio},
    tem ao menos um claro e um escuro entre eles.
  \item $e$ possui tr\^{e}s vizinhos $\{b,c,d\}$ j\'{a} coloridos e, analogamente,
    tem ao menos um claro e um escuro.
  \item $b$, $c$ e $d$ possu\'{i}am quatro vizinhos em algum $K_5$ anterior e, pela hip\'{o}tese indutiva,
    j\'{a} t\^{e}m ambas as cores.
\end{itemize}

Portanto, \textbf{independentemente da cor atribu\'{i}da \`{a} nova aresta $ae$}, todos os cinco
v\'{e}rtices do novo $K_5$ continuam tendo ao menos um vizinho claro e um vizinho escuro.
A Proposi\c{c}\~{a}o~\ref{prop:equilibrio} \'{e} preservada no Caso U1.

\medskip
\begin{center}
\begin{tikzpicture}[scale=1.1,
  vx/.style={circle,draw,fill=white,minimum size=14pt,font=\tiny\strut}]
  \node[vx] (a) at (90:1.4)   {$a$};
  \node[vx] (b) at (18:1.4)   {$b$};
  \node[vx] (c) at (-54:1.4)  {$c$};
  \node[vx] (d) at (-126:1.4) {$d$};
  \node[vx] (e) at (162:1.4)  {$e$};
  % 9 solid edges of K5^- (without a-e)
  % Red edges (5 of the 9)
  \draw[red,  line width=1.8pt] (a)--(b) (b)--(c) (c)--(d) (d)--(e);
  % Blue edges (4 of the 9)
  \draw[blue, line width=1.8pt] (a)--(c) (a)--(d) (b)--(d) (b)--(e) (c)--(e);
  % Missing edge a-e (dashed, being added)
  \draw[dashed, line width=2pt, black!60] (a)--(e);
  \node[font=\tiny, left=2pt] at ($(a)!0.5!(e)$) {$ae$};
\end{tikzpicture}
\end{center}

\noindent
A aresta tracejada $ae$ \'{e} a \'{u}nica aresta adicionada em U1.
Qualquer cor (claro ou escuro) pode ser atribu\'{i}da a $ae$;
a propriedade de equil\'{i}brio \'{e} mantida pois $a$ e $e$
j\'{a} possu\'{i}am ambas as cores entre seus vizinhos $\{b,c,d\}$.

\medskip
\noindent
Conclu\'{i}mos que, ap\'{o}s todas as etapas de expans\~{a}o ($E_0$, $E_1$, $E_2$, $U_1$),
a colora\c{c}\~{a}o do bloco $B$ satisfaz a Proposi\c{c}\~{a}o~\ref{prop:equilibrio},
e portanto, pela Observa\c{c}\~{a}o~\ref{rem:equilibrio-suficiente},
evita $K_3$ monocrom\'{a}tico, $K_4$ arco-\'{i}ris e $K_5$ lexicogr\'{a}fico.


\end{proof}




\vspace{1cm}

% ============================================================
\printbibliography
% ============================================================

\end{document}
